%---------- Inleiding ---------------------------------------------------------

\section{Inleiding}%
\label{sec:inleiding}

Bij het zoeken naar mogelijke aanwinsten moeten scouts bepalen welke spelers aansluiten bij een gewenst spelersprofiel. Deze bachelorproef vertrekt vanuit de vraag hoe beschikbare wedstrijd- en spelersstatistieken kunnen worden ingezet om die zoektocht te ondersteunen. De focus ligt daarbij niet uitsluitend op het weergeven van cijfers, maar op het vertalen van deze gegevens naar gerichte zoekmogelijkheden en interpreteerbare spelersvergelijkingen.

Dit onderzoek richt zich op het ontwerpen, realiseren en evalueren van een datagedreven scoutingomgeving voor de BNXT League. Wedstrijd- en spelersgegevens worden centraal verzameld, opgeschoond en gemodelleerd. Een interactief Power BI-dashboard maakt deze gegevens toegankelijk via spelersoverzichten, filters, detailpagina's en vergelijkingsschermen. De primaire doelgroep bestaat uit scouts die spelers willen identificeren voor verdere beoordeling.

Naast deze algemene analysemogelijkheden staan twee scoutingfunctionaliteiten centraal. De eerste ondersteunt het zoeken naar spelers met een gelijkaardig statistisch profiel als een geselecteerde referentiespeler. De tweede ondersteunt het signaleren van veelbelovende spelers op basis van expliciete prestatiecriteria. Hiervoor worden machine learning-technieken en aanvullende analysemethoden onderzocht.

Een gelijkaardig statistisch profiel verwijst binnen dit onderzoek naar overeenkomsten in geselecteerde prestatie-indicatoren. Dit betekent niet automatisch dat spelers dezelfde speelstijl, kwaliteit of geschiktheid voor een team hebben. Veelbelovende spelers worden beschouwd als spelers die op basis van onderbouwde criteria in aanmerking komen voor verdere scouting. Deze omschrijving is niet beperkt tot jonge spelers en houdt geen voorspelling van toekomstig spelerssucces in.

De centrale onderzoeksvraag luidt als volgt:

\begin{quote}
Hoofdonderzoeksvraag: Hoe kan een datagedreven scoutingomgeving voor de BNXT League worden ontworpen, gerealiseerd en geëvalueerd om scouts te ondersteunen bij het identificeren van spelers met een gelijkaardig statistisch profiel en het signaleren van veelbelovende spelers?
\end{quote}

Om deze hoofdonderzoeksvraag te beantwoorden, worden het probleemdomein en het oplossingsdomein opgesplitst in de volgende deelvragen.

\subsection{Deelvragen: probleemdomein}

\begin{enumerate}
  \item Hoe beoordelen scouts mogelijke aanwinsten en welke informatie hebben zij nodig bij het zoeken naar gelijkaardige en veelbelovende spelers?
  \item Welke gegevens over wedstrijden en spelers zijn beschikbaar voor de BNXT League en welke beperkingen of kwaliteitsproblemen bevatten deze gegevens?
  \item Welke prestatie-indicatoren en criteria zijn relevant om statistische spelersgelijkenis en veelbelovende spelers binnen dit onderzoek af te bakenen?
\end{enumerate}

\subsection{Deelvragen: oplossingsdomein}

\begin{enumerate}
  \item Hoe kan een data-architectuur en pipeline worden ontworpen om BNXT-wedstrijd- en spelersstatistieken centraal te verzamelen, op te schonen en te modelleren?
  \item Welke machine learning-technieken en aanvullende analysemethoden zijn geschikt om gelijkaardige spelers te identificeren en veelbelovende spelers te signaleren?
  \item Welke visualisaties, filters en toelichtingen ondersteunen scouts bij het interpreteren en vergelijken van de voorgestelde kandidaten?
\end{enumerate}

Het onderzoek levert een proof-of-concept op waarin een centrale datalaag en een interactief dashboard voor algemene spelersanalyse samenkomen met de twee scoutingfunctionaliteiten. De omgeving ondersteunt de voorselectie van kandidaten en laat scouts toe om de voorgestelde spelers verder te onderzoeken. De gebruikte gegevens, indicatoren en analysemethoden worden gedocumenteerd, zodat de resultaten kunnen worden geïnterpreteerd en gereproduceerd. De uiteindelijke spelersbeoordeling blijft bij de scout.

\section{Literatuurstudie}%
\label{sec:literatuurstudie}

Datagedreven analyses bieden aanknopingspunten om spelersbeoordeling te ondersteunen. Demenius en Kreivytė onderzochten de opvattingen van dertig managers, coaches en assistent-coaches van tien Litouwse mannenbasketbalteams. De respondenten stonden overwegend positief tegenover geavanceerde basketbalanalyses en zagen mogelijkheden om tegenstanders en mogelijke nieuwe spelers beter te beoordelen en financiële middelen gerichter in te zetten \autocite{demenius2017benefits}. Deze bevindingen ondersteunen de relevantie van analytische informatie, maar vormen geen rechtstreeks bewijs dat een specifiek dashboard betere rekruteringsbeslissingen oplevert.

Rekrutering omvat daarnaast meer dan wedstrijdstatistieken. Banucha en Janasz beschrijven in een casestudy bij een Poolse professionele basketbalclub hoe prestatiegegevens, psychologische kenmerken, scoutingnetwerken en financiële en sportieve doelstellingen samenkomen bij de selectie van spelers en coaches \autocite{banucha2025strategic}. Deze bredere context is relevant voor de afbakening van de bachelorproef: de scoutingomgeving ondersteunt de statistische voorselectie, maar omvat niet alle aspecten van een rekruteringsbeslissing.

Onderzoek naar prestatieprofielen biedt een vertrekpunt voor de keuze van indicatoren en analysemethoden. Dehesa et al. onderzochten wedstrijdgebonden profielen van 472 NBA-spelers in 535 wedstrijden met een eindverschil van maximaal acht punten. De onderzochte variabelen omvatten onder meer speeltijd, wedstrijdtempo, offensive en defensive rating, effective field-goal percentage en turnover percentage. Met een clusteranalyse werden vijf profielen tijdens de reguliere competitie en vier tijdens de play-offs onderscheiden \autocite{dehesa2019key}. Deze studie biedt een voorbeeld van statistische profielvorming, maar de gebruikte indicatoren worden niet zonder verdere beoordeling overgenomen voor de BNXT League.

Een andere toepassing van spelersgelijkenis wordt beschreven door Muniz en Flamand. Zij combineren k-meansclustering met een gewogen netwerk en de Louvainmethode om spelersprofielen in de NBA te onderscheiden. Spelers vormen de knopen van het netwerk. De gewichten tussen spelers worden bepaald door de mate waarin zij samen in clusters terechtkomen voor verschillende categorieën basketbalgegevens. De methode wordt toegepast op zes NBA-seizoenen en maakt onder meer gebruik van trackinggegevens \autocite{muniz2022weighted}.

Deze studie biedt een aanknopingspunt voor het onderzoeken van spelersgelijkenis op basis van meerdere statistische dimensies. De aanpak wordt echter niet rechtstreeks overgenomen voor de BNXT League, aangezien de beschikbare gegevens kunnen verschillen. Binnen deze bachelorproef wordt onderzocht welke profielen met de beschikbare statistieken kunnen worden samengesteld en hoe deze scouts kunnen ondersteunen bij het identificeren van kandidaten.

Voor talentsignalering is ook de afbakening van de beschikbare informatie belangrijk. Han et al. onderzochten in een systematische review en meta-analyse achttien studies over antropometrische, fysiologische en fysieke prestatiekenmerken bij talentidentificatie in basketbal. De review rapporteert verschillen tussen onderzochte spelersgroepen voor onder meer lengte, lichaamsmassa, wendbaarheid, snelheid en sprongprestaties \autocite{han2023basketball}. Deze kenmerken zijn niet allemaal beschikbaar in gewone wedstrijdstatistieken.

Daarom wordt binnen deze bachelorproef onderscheid gemaakt tussen volledige talentbeoordeling en het signaleren van spelers op basis van beschikbare statistische gegevens. Het onderzoek richt zich op deze tweede toepassing. Veelbelovend wordt vertaald naar expliciete criteria die worden besproken met een werkveldexpert. De voorgestelde kandidaten worden beschouwd als spelers voor verdere opvolging, niet als bewezen toekomstige succesvolle spelers.

Op basis van de literatuur worden prestatie-indicatoren, mogelijke analysemethoden en evaluatiecriteria geselecteerd. De bijdrage van de bachelorproef bestaat uit het ontwerpen en evalueren van een concrete scoutingomgeving voor de BNXT League.

%---------- Methodologie ------------------------------------------------------
\section{Methodologie}%
\label{sec:methodologie}

Het onderzoek wordt opgezet als een ontwerpgerichte en deels vergelijkende studie waarin een datagedreven scoutingomgeving voor de BNXT League wordt ontworpen, gerealiseerd en geëvalueerd. De methode combineert een literatuurstudie, requirementsanalyse, gegevensverwerking, modelontwikkeling en de implementatie van een Power BI-dashboard.

\subsection{Fase 1: Literatuurstudie en uitwerking deelvragen}

In de eerste fase wordt onderzocht welke prestatie-indicatoren en analysemethoden relevant zijn voor spelersvergelijking en talentsignalering. Daarnaast wordt een werkveldexpert betrokken om de informatiebehoeften van scouts te bespreken en de twee scoutingtaken concreet af te bakenen.

Daarbij wordt bepaald wat binnen het onderzoek wordt verstaan onder een gelijkaardig statistisch profiel en een veelbelovende speler. Ook worden de evaluatiecriteria en een eenvoudige referentiemethode vastgelegd.

De belangrijkste resultaten van deze fase zijn:
\begin{itemize}
  \item Een overzicht van relevante en mogelijk beschikbare prestatie-indicatoren.
  \item Functionele requirements voor de scoutingomgeving.
  \item Twee afgebakende scoutingcases en bijbehorende beoordelingscriteria.
  \item Een beschrijving van de referentiemethode waarmee de ontwikkelde aanpak wordt vergeleken.
\end{itemize}

\subsection{Fase 2: Dataverzameling, opschoning en modellering}

In de tweede fase worden historische wedstrijd- en spelersgegevens uit de BNXT League verzameld uit geselecteerde publieke of aangeleverde bronnen. De onderzochte seizoenen, gebruikte bronnen en dekking van de dataset worden gedocumenteerd.

Python en pandas worden gebruikt voor het opschonen en transformeren van de gegevens. Daarbij worden onder meer ontbrekende waarden, dubbele records, inconsistente identificaties en teamwissels onderzocht. SQL Server wordt gebruikt om de gegevens centraal op te slaan en te modelleren.

Afgeleide indicatoren worden voorzien van een definitie, berekeningswijze en beschrijving van hun beperkingen. Er wordt onderzocht hoe verschillen in speeltijd worden verwerkt en welke schaaltransformaties nodig zijn voor de analysemethoden. Minimumvoorwaarden voor opname in de analyses, zoals gespeelde minuten of wedstrijden, worden verantwoord.

Exacte individuele balbezitten worden niet verondersteld beschikbaar te zijn. Indicatoren die aanvullende gegevens vereisen, worden alleen opgenomen wanneer de benodigde gegevens voldoende betrouwbaar beschikbaar zijn.

Deze fase resulteert in een gedocumenteerde dataset, een datamodel en reproduceerbare verwerkingsscripts.

\subsection{Fase 3: Ontwerp en implementatie van de scoutingomgeving}

In de derde fase worden de analysemethoden ontwikkeld en geïntegreerd in een Power BI-dashboard bovenop de centrale datalaag.

Voor het identificeren van gelijkaardige spelers wordt een nearest-neighborbenadering onderzocht. Hiermee worden kandidaten gerangschikt ten opzichte van een geselecteerde referentiespeler. De invloed van de gekozen indicatoren, schaaltransformaties, afstandsmaat en vergelijkingsgroep wordt onderzocht.

Daarnaast wordt clustering onderzocht om statistische spelersprofielen te onderscheiden. Er wordt nagegaan of deze profielen een bruikbare bijdrage leveren aan spelersvergelijking en talentsignalering. Voor het signaleren van veelbelovende spelers worden expliciete prestatiecriteria opgesteld. Daarbij wordt duidelijk onderscheiden welke resultaten voortkomen uit machine learning en welke uit vaste filters of selectieregels.

De analyses worden uitgevoerd in Python. SQL Server ondersteunt de centrale datalaag. Power BI wordt gebruikt voor het semantische model, DAX-measures en visualisaties.

Op basis van de requirements worden de volgende onderdelen voorzien:
\begin{itemize}
  \item Een spelersoverzicht met filters voor relevante kenmerken, zoals seizoen, team, positie, leeftijd en speeltijd.
  \item Detailpagina's met kernstatistieken en prestatie-indicatoren.
  \item Vergelijkingsschermen waarin geselecteerde spelers naast elkaar worden geplaatst.
  \item Een zoekfunctionaliteit voor spelers met een gelijkaardig statistisch profiel.
  \item Een overzicht van spelers die voldoen aan de criteria voor verdere scouting.
\end{itemize}

Bij voorgestelde kandidaten worden de onderliggende indicatoren, overeenkomsten, verschillen en beschikbare minuten en wedstrijden getoond. De gebruikte vergelijkingsgroep en selectievoorwaarden blijven zichtbaar.

Indien de gegevens voldoende dekking bieden, worden periode- of seizoensvergelijkingen toegevoegd als ondersteunende context. Deze vormen geen afzonderlijk voorspellend model.

\subsection{Fase 4: Evaluatie en gebruiksbeoordeling}

In de vierde en laatste fase wordt het prototype technisch en inhoudelijk geëvalueerd.

De technische evaluatie omvat:
\begin{itemize}
  \item Controles van de datakwaliteit en gegevenskoppelingen.
  \item Verificatie van berekende indicatoren aan de hand van geselecteerde voorbeelden.
  \item Controle van filters, rangschikkingen en dashboardinteracties.
  \item Beoordeling van de analysemethoden op passende technische kwaliteitsmaten en gevoeligheid voor gewijzigde instellingen.
  \item Meting van relevante verwerkings- en responstijden.
\end{itemize}

De ontwikkelde aanpak wordt vergeleken met een eenvoudige referentiemethode op basis van vooraf vastgelegde filters en rangschikkingen. Hiermee wordt onderzocht welke bijdrage de analysemethoden leveren bovenop een eenvoudige selectie.

De gebruiksbeoordeling gebeurt met een werkveldexpert aan de hand van twee vooraf gedefinieerde scoutingcases:
\begin{enumerate}
  \item Identificeer spelers met een gelijkaardig statistisch profiel als een gekozen referentiespeler.
  \item Identificeer spelers voor verdere scouting op basis van een vooraf vastgelegd zoekprofiel.
\end{enumerate}

De beoordeling richt zich op de relevantie van de voorgestelde kandidaten, de begrijpelijkheid van de toelichting en de bruikbaarheid van de omgeving. Waar haalbaar worden taakduur en het aantal relevant beoordeelde kandidaten geregistreerd.

De resultaten worden beschreven met vermelding van de onderzochte cases, het aantal beoordelaars en de beperkingen van de dataset. De evaluatie vormt geen bewijs van toekomstig spelerssucces of algemene bruikbaarheid voor alle scoutingopdrachten.

\subsection{Tools}

Het onderzoek maakt gebruik van Python, pandas en scikit-learn voor gegevensverwerking en modelontwikkeling. SQL Server wordt gebruikt voor de centrale opslag en modellering. Power BI ondersteunt de interactieve rapportering. GitHub wordt gebruikt voor versiebeheer van scripts, experimenten en documentatie.

%---------- Verwachte resultaten ----------------------------------------------
\section{Verwachte resultaten en conclusie}%
\label{sec:verwachte_resultaten}

Het verwachte resultaat is een proof-of-concept van een datagedreven scoutingomgeving voor de BNXT League. Het Power BI-dashboard maakt het mogelijk om spelers te filteren op relevante kenmerken, statistieken te analyseren en spelersprofielen naast elkaar te plaatsen.

Daarnaast ondersteunt de omgeving twee gerichte scoutingtaken: het identificeren van spelers met een gelijkaardig statistisch profiel en het signaleren van veelbelovende spelers voor verdere opvolging. Bij de voorgestelde kandidaten worden de onderliggende indicatoren en de omvang van de beschikbare gegevens inzichtelijk gemaakt.

Naast het prototype levert het onderzoek documentatie van de dataset, het datamodel, de indicatoren en de analysemethoden op. De technische tests en gebruiksbeoordeling beschrijven in welke mate de omgeving consistente, interpreteerbare en bruikbare resultaten oplevert binnen de onderzochte scoutingcases.

De beoogde meerwaarde voor scouts ligt in het gericht doorzoeken van spelersgegevens en het onderbouwen van een voorselectie. Of de ontwikkelde methoden daadwerkelijk meerwaarde bieden ten opzichte van de eenvoudige referentiemethode, wordt op basis van de evaluatie beoordeeld.

De conclusies worden begrensd door de beschikbare gegevens, de gekozen scoutingcases en de omvang van de werkveldbeoordeling. De omgeving ondersteunt de expertise van scouts, maar vervangt geen volledige spelersbeoordeling en garandeert geen toekomstig spelerssucces.